\clearpage\pagestyle{otMain}\setlength{\otBlocEspacement}{2.5mm}\fontsize{11}{15}\selectfont\corrigesChapitre{2}{Nombres rationnels et fractions}
\begin{blocLivre}{cor-2-1}\begin{corrige}{1}{Des fractions égales}\demarche{Agir de la même façon sur les deux termes.}
$42$ et $56$ sont dans la table de 14 : $\frac{42}{56}=\res{\frac34}$. De même $\frac{36}{-48}=\frac{3}{-4}=\res{-\frac34}$. Enfin on multiplie les deux termes de $\frac{-3}4$ par 5 : $\res{\frac{-15}{20}}$.
\end{corrige}\end{blocLivre}
\begin{blocLivre}{cor-2-2}\begin{corrige}{2}{Entier, décimal ou non}\demarche{Effectuer la division du numérateur par le dénominateur.}
$\frac{12}4=3$ et $\frac{-18}6=-3$ sont $\res{\text{entiers}}$, donc aussi décimaux. $\frac7{20}=\res{0{,}35}$ est décimal non entier. $\frac13=0{,}333\ldots$ n’est $\res{\text{ni entier ni décimal}}$ : la division ne s’arrête jamais.
\end{corrige}\end{blocLivre}
\begin{blocLivre}{cor-2-3}\begin{corrige}{3}{L’opposé d’un quotient}\demarche{Utiliser la définition de l’opposé.}
$\frac57+\frac{-5}7=\frac07=\res0$ : les deux nombres sont opposés, donc $\frac{-5}7=-\frac57$. En multipliant les deux termes par $-1$ : $\frac3{-8}=\res{\frac{-3}8=-\frac38}$.
\end{corrige}\end{blocLivre}
\begin{blocLivre}{cor-2-4}\begin{corrige}{4}{Additionner}\demarche{Réduire au même dénominateur.}
Le dénominateur commun 12 donne $\frac9{12}-\frac{10}{12}=\res{-\frac1{12}}$. Le résultat négatif est cohérent puisque $\frac34<\frac56$.
\end{corrige}\end{blocLivre}
\begin{blocLivre}{cor-2-5}\begin{corrige}{5}{Multiplier}\demarche{Décomposer pour simplifier avant de multiplier.}
Le résultat est négatif. On simplifie avant de multiplier : $\frac{4\times21}{7\times10}=\frac{4\times3\times7}{7\times2\times5}=\frac65$. Donc le produit vaut $\res{-\frac65}$.
\end{corrige}\end{blocLivre}
\begin{blocLivre}{cor-2-6}\begin{corrige}{6}{Une fraction d’une fraction}\demarche{Représenter puis multiplier les fractions.}
Les $\frac34$ du rectangle sont 3 colonnes, soit 9 cases sur 12. Les $\frac23$ de cette partie sont 2 lignes de ces 3 colonnes, soit 6 cases. Le calcul est $\frac23\times\frac34=\frac6{12}=\res{\frac12}$ du rectangle.
\label{fig-cor-2-6}\figureGeo[corrige]{fraction-rectangle}
\end{corrige}\end{blocLivre}
\begin{blocLivre}{cor-2-7}\begin{corrige}{7}{Trouver un inverse}\demarche{Transformer la division en multiplication.}
$\left(-\frac53\right)\times\left(-\frac35\right)=1$, donc l’inverse est $\res{-\frac35}$. Le quotient vaut $\frac23\times\left(-\frac35\right)=\res{-\frac25}$. Seul le diviseur est remplacé par son inverse.
\end{corrige}\end{blocLivre}
\begin{blocLivre}{cor-2-8}\begin{corrige}{8}{Opposé ou inverse}\demarche{Somme nulle pour l’opposé, produit égal à 1 pour l’inverse.}
Opposés : $\res{-4}$, $\res{\frac16}$ et $\res{-0{,}25}$, car chaque somme vaut 0. Inverses : $\res{\frac14}$, $\res{-6}$ et $\res4$, car chaque produit vaut 1. L’inverse d’un nombre négatif est négatif.
\end{corrige}\end{blocLivre}
\begin{blocLivre}{cor-2-9}\begin{corrige}{9}{Un calcul enchaîné}\demarche{Respecter la priorité du produit.}
La multiplication est prioritaire : $\frac23\times\frac94=\frac{18}{12}=\frac32$. Donc $\frac56-\frac32=\frac56-\frac96=\res{-\frac23}$.
\end{corrige}\end{blocLivre}
\begin{blocLivre}{cor-2-10}\begin{corrige}{10}{Diviser une somme}\demarche{Terminer la parenthèse avant d’inverser le diviseur.}
La parenthèse vaut $\frac14$. On multiplie ensuite par l’inverse $\frac85$ : $\frac14\times\frac85=\res{\frac25}$. On peut vérifier que $\frac25\times\frac58=\frac14$.
\end{corrige}\end{blocLivre}
\begin{blocLivre}{cor-2-11}\begin{corrige}{11}{Qui a bu le plus ?}\demarche{Exprimer chaque part comme fraction de la bouteille entière.}
Après Inès, il reste $1-\frac25=\frac35$ de la bouteille. Hugo boit $\frac23\times\frac35=\frac25$ de la bouteille : $\res{\text{ils ont bu autant}}$. Il reste $1-\frac25-\frac25=\res{\frac15}$ de la bouteille. La fraction de Hugo porte sur le reste, pas sur la bouteille entière.
\end{corrige}\end{blocLivre}
\begin{blocLivre}{cor-2-12}\begin{corrige}{12}{Lire une portion du reste}\demarche{Préciser à quelle quantité s’applique chaque fraction.}
Après le premier usage, il reste $\left(1-\frac25\right)\times60=36$ L. On utilise ensuite $\frac13\times36=12$ L, donc il reste $\res{24\ \mathrm L}$. La seconde fraction ne porte pas sur les 60 L initiaux.
\end{corrige}\end{blocLivre}
\begin{blocLivre}{cor-2-13}\begin{corrige}{13}{Retrouver une quantité}\demarche{Diviser par la fraction connue.}
Si $B$ est le budget, $\frac7{12}\times B=175$. Donc $B=175:\frac7{12}=175\times\frac{12}7=\res{300}$ euros. La part restante vaut $300-175=\res{125}$ euros. Contrôle : $\frac7{12}\times300=175$.
\end{corrige}\end{blocLivre}
\begin{blocLivre}{cor-2-14}\begin{corrige}{14}{Une affirmation à discuter}\demarche{Tester une fraction comprise entre 0 et 1.}
$6:\frac12=6\times2=12>6$ est un $\res{\text{contre-exemple}}$. Pour un nombre positif, diviser par un nombre strictement compris entre 0 et 1 revient à multiplier par son inverse, supérieur à 1 : le quotient augmente. Diviser par un nombre supérieur à 1 le diminue.
\end{corrige}\end{blocLivre}
\begin{blocLivre}{cor-2-15}\begin{corrige}{15}{Trois parts successives}\demarche{Exprimer le reste final comme fraction du total.}
Après Lina, il reste $\frac34$ du total. Après Sami, il reste $\frac13$ de ce reste, soit $\frac13\times\frac34=\frac14$ du total. Donc le total vaut $36\times4=\res{144}$ euros. Lina reçoit $\res{36}$ euros, Sami $\frac23\times108=\res{72}$ euros. Contrôle : $36+72+36=144$.
\end{corrige}\end{blocLivre}
